% FILE: 02_lineage_and_context.tex
% PROJECT: crosswalks
% ROLE: type-law-crosswalk-between-standards-and-execution

\section{Lineage and Context}
\label{sec:crosswalks-lineage}

\subsection{2016 Language-Model Lesson}
\label{sec:crosswalks-lineage-2016}

The thesis lineage begins with a 2016 language-model experiment whose conclusion was that local
text imitation does not conserve consequence. A model that predicts the next token in a sequence
reproduces the surface distribution of a process without preserving the causal structure that
makes the process lawful. The model can emit output that looks like a valid event trace while
the underlying constraints --- precedence, response, coexistence --- are violated.

The crosswalks project is a formal elaboration of this lesson applied to the translation layer
of process intelligence systems. When \textsc{PM4Py} converts an \textsc{OCEL} log to \textsc{XES}
by silently dropping non-chosen object types, it performs the same operation structurally: it
produces an artifact that looks like a complete event log while the causal links to non-chosen
objects are absent. The \textsc{XES} output imitates the surface of a case-centric log without
conserving the object-centric consequence. The crosswalk \texttt{OCEL\_TO\_XES.md} names this
as information destruction and requires a \texttt{LossReport} to accompany it.

The 2016 lesson is thus not merely about language models. It is about the general danger of
format-preserving projections that do not conserve the structural properties that make the source
artifact evidentially meaningful.

\subsection{Enterprise Process Gap}
\label{sec:crosswalks-lineage-enterprise}

The enterprise process gap that motivates this research is the absence of a type-safe translation
layer between heterogeneous process formalisms. Enterprise process intelligence deployments
routinely combine data from systems that export in \textsc{BPMN}, \textsc{XES}, \textsc{OCEL},
and Declare formats. When these formats are combined, translation is unavoidable. The question is
whether the translation is audited.

The crosswalks project addresses this gap by defining, for each translation direction, the exact
categories of loss that a production system must account for. The \texttt{BPMN\_TO\_WFNET.md}
crosswalk lists seven categories of hard drop (annotations, data objects, message flows, lane
structure, gateway conditions, event types) and two categories of structural approximation
(\textsc{OR}-gateway over-approximation, loop marker restructuring). An enterprise system that
ingests \textsc{BPMN} and exports conformance scores must account for all nine categories or
declare, under a named \texttt{LossPolicy}, which it permits.

\textbf{[AUTHOR THESIS]:} The enterprise process gap is not a data integration problem. It is
a type-law problem. Closing it requires not better ETL pipelines but formal projection contracts
that make loss refusable, nameable, and auditable.

\subsection{Relationship to the Chatman Equation}
\label{sec:crosswalks-lineage-chatman}

The Chatman Equation $A = \mu(O^*)$ states that an autonomous agent $A$ is a function of the
object model $O^*$ it observes. The crosswalks project grounds the observational surface of
$O^*$: it defines which structural properties of a process are preserved when $O^*$ is projected
from one formalism to another.

\textbf{[INTERPRETATION]:} If the projection from \textsc{OCEL} to \textsc{XES} silently drops
object types, then the $O^*$ that the agent observes is an approximation of the true object
model. The conformance verdicts the agent emits are verdicts about the approximation. The
Chatman Equation is satisfied only if $O^*$ is faithfully represented --- which requires the
projection to be lawful under an explicit \texttt{LossPolicy}. The crosswalks project defines
what ``faithful'' means for each formalism pair.

\subsection{Relationship to Receipts and Replay}
\label{sec:crosswalks-lineage-receipts}

The receipt and replay infrastructure of the process intelligence foundry requires that every
evidence artifact carry a provenance chain: a \textsc{BLAKE3} hash over the serialised payload,
state, witness, and epoch, verified by an Ed25519 signature (source: \texttt{alive\_001\_to\_iso\_compliance.md},
\textsection 2.3, citing \textsc{ISO/IEC} 27001). This cryptographic binding is only meaningful
if the artifact it binds faithfully represents its source.

When a \textsc{BPMN} process with \textsc{OR}-gateways is projected to a WF-net by silent
approximation, the receipt chain binds a structurally compromised artifact. Replay against this
artifact produces conformance scores that cannot be attributed to the original \textsc{BPMN}
process. The crosswalks project defines the admission gate that must be passed before a projected
artifact enters the receipt chain: the \texttt{LossPolicy} must be declared, the \texttt{LossReport}
must be produced, and the projection must be admitted under the typed witness before its receipt
is computed.

The \texttt{OCEL\_TO\_XES.md} crosswalk formalises this as a five-step covenant: raw bytes
\(\to\) \texttt{FormatEnvelope} \(\to\) \texttt{Admission<OcelLog, Ocel20>} \(\to\) lossy
\texttt{Project} with named \texttt{LossPolicy} \(\to\) \texttt{LossReport} emitted \(\to\)
projected \texttt{XesLog} admitted for export. No step can be skipped.
